\section{Introduction}
\label{sec:intro}

Compact heat exchangers, micro-reactors and electronics coolers are increasingly built with polymer
walls, which combine low thermal conductivity with a high coefficient of thermal expansion. For the
same imposed heat flux, such a wall sustains a through-thickness temperature difference that is
far larger than that of a metallic wall of equal thickness, and linear
thermoelasticity associates this difference with a free thermal curvature of the wall
\citep{boley1960theory,timoshenko1970elasticity}. When the wall is as thick as, or thicker than, the
channel is high, the curvature is not a small perturbation of the passage: it modifies the flow
cross-section, and heat-transfer predictions that assume a fixed geometry become geometry-dependent
in a thermally self-consistent way.

Fluid--structure interaction (FSI) with heat transfer has mostly been studied for deformation
driven by the fluid. \citet{soti2015steady} computed laminar channel flow past a flexible thin plate
and showed that the flow-induced motion passively raises the heat transfer; \citet{alsabery2019numerical}
and \citet{saleh2025laminar} treated transient natural convection in cavities with flexible walls;
\citet{christov2018flow} derived the steady relation between flow rate and pressure drop in a
deformable shallow channel, again with pressure as the only load. The structure is thermally passive
in these studies. Thermally driven structural motion has been computed for heated flexible
cantilevers in a differentially heated cavity \citep{kumar2026thermo}, a confined natural-convection
configuration rather than a forced-convection channel. Conjugate heat transfer in rigid thick-wall channels, finally,
has been studied \citep{rahimi2024conjugate}, but there the wall does not move.

The configuration studied here differs in the mechanism and in the flow class. A finite segment of
one wall of a laminar forced-convection slit is a thin thermoelastic strip, heated on its outer
face. Conduction across the strip sets up a temperature difference; the associated eigen-curvature,
acting on a strip with pinned ends, bows it \emph{away} from the passage; the widened passage
changes the near-wall velocity and temperature gradients and therefore the interfacial heat flux,
which is one of the boundary data of the conduction problem in the strip. Within the literature
examined for this work we did not find a study that combines (i) the thermal eigen-curvature of the
wall as the principal deformation mechanism, (ii) a closed steady loop through the conjugate
interface, and (iii) a confined forced-convection channel. This positioning is deliberately bounded:
the literature search was not exhaustive, and it is not a claim of priority.

The paper makes five contributions, and states clearly what it does not do.
\begin{enumerate}
\item A complete steady formulation with one consistent set of sign and interface conventions
(Section~\ref{sec:formulation}).
\item A reduced-order thermoelastic analysis of the strip in closed form: free curvature,
pinned bowing, thermal membrane force and buckling, the beam--column amplification of the bow, and
the ratio of thermal to hydraulic deflection (Section~\ref{sec:reduced}).
\item A baseline operating point evaluated from first principles, with the heat-flux ceilings
imposed by deformation, stress and buckling, and an \emph{indicative} estimate of what the
deformation does to the flow (Section~\ref{sec:baseline}).
\item A specified finite-element formulation and fixed-point algorithm for the coupled problem, with
the treatment of the interface variables made explicit (Section~\ref{sec:fe}).
\item A plane-strain Q8 verification of the wall idealisation: a mesh-sensitivity study of the free bow at a
fixed sampling height, and an axially restrained, geometrically nonlinear solution of the baseline wall
(Section~\ref{sec:q8}).
\end{enumerate}
\textbf{No coupled fluid--structure computation is reported.} The coupled solver is specified, not
executed; every statement about the flow in Section~\ref{sec:baseline} is a reduced-order estimate
and is labelled as such. The paper is organised accordingly: formulation, analysis, baseline, finite
element method and verification, and limitations.
